\chapter{现代多核心 CPU 的执行结构}

现代 CPU 不是“一个更快的核心”那么简单。它包含多个核心、每个核心内部的流水线和执行端口、多级缓存、共享末级缓存、内存控制器、硬件预取器、TLB、互连结构和可能存在的 NUMA 节点。推理引擎运行在这样的机器上，性能取决于所有这些结构如何共同工作。

\section{核心、线程和拓扑}

一个物理核心通常包含前端取指和解码、乱序执行窗口、整数和浮点执行单元、加载和存储单元、寄存器文件以及私有 L1/L2 缓存。超线程或 SMT 让一个物理核心暴露为多个逻辑 CPU，多个硬件线程共享部分核心资源。对计算密集型内核而言，SMT 不一定带来线性加速；对等待内存的内核而言，SMT 有时能隐藏一部分延迟。

Linux 把可调度对象放到逻辑 CPU 上运行。读者在 benchmark 时看到的“线程数”，不等于一定占用了同样数量的物理核心。若线程在不同核心之间迁移，缓存热度会被破坏；若多个重计算线程落在同一物理核心的两个 SMT 线程上，浮点执行资源会竞争。因此推理引擎 benchmark 要记录核心拓扑，并在必要时使用线程亲和性。

\begin{mentalmodel}
多核心 CPU 是一组共享内存的执行工人。每个工人有自己的近身工具箱，也共享仓库和道路。任务切分好不好，不只看工人数，还要看工具箱是否够用、道路是否拥堵、工人是否反复换位置。
\end{mentalmodel}

\section{流水线和乱序执行}

CPU 不会逐行机械执行 C++ 代码。前端把指令取出、解码成内部微操作，后端在依赖允许时乱序调度到不同执行端口。若两条指令没有数据依赖，硬件可以并行推进；若一条指令依赖上一条结果，就形成依赖链。点积中的单累加器就是典型依赖链，多个累加器能让硬件看到更多独立工作。

推理引擎中的算子常常有大量重复循环。循环体是否容易解码，地址计算是否复杂，加载是否及时返回，FMA 是否有足够独立输入，都会影响单核吞吐。编译器优化和手写内核的目的，不是让代码看起来短，而是让硬件后端有连续、规则、独立的工作可以调度。

\section{共享缓存和一致性}

每个核心通常有私有 L1 数据缓存和 L2 缓存，多个核心共享 L3 或 LLC。核心之间通过缓存一致性协议维护同一内存地址的可见性。读共享数据通常比较便宜，多个核心同时写相邻数据就可能触发 cache line 来回转移。false sharing 正是这种现象：逻辑上不同变量，物理上落在同一 cache line，导致一致性流量。

推理引擎中，权重通常是只读共享数据，适合被多个线程读取；输出张量应该按块分配给不同线程写，避免写同一 cache line；线程本地统计量要考虑对齐和填充。多线程性能问题往往不是“锁慢”这么简单，而是共享写、同步点和缓存一致性共同作用。

\section{内存控制器和 NUMA}

单 socket 机器通常有一个或多个内存通道，多 socket 机器可能有 NUMA 节点。线程访问本地节点内存和远端节点内存的延迟、带宽不同。大模型权重和 KV Cache 占用很大，NUMA 放置会影响吞吐和尾延迟。即使第一阶段 LCQI 模型很小，本册也要提前建立 NUMA 意识，因为后续扩展到大模型时这个问题无法回避。

\begin{warningbox}
不要把“核心数”当成唯一性能指标。同样 16 个逻辑 CPU，可能是 8 个物理核心加 SMT，也可能是大小核混合，也可能跨 NUMA 节点。性能报告必须写清楚拓扑，否则线程扩展结果无法解释。
\end{warningbox}

\section{和推理引擎的关系}

LCQI 的每个模块都会受到 CPU 结构影响。模型加载主要受 IO 和解析影响；张量遍历受缓存和 TLB 影响；量化线性层受内存带宽、SIMD 和依赖链影响；线程池受调度、同步和缓存一致性影响；benchmark 受频率、亲和性和系统噪声影响。理解 CPU 结构，是为了把这些现象放回正确层级。

\begin{exercisebox}
在本地 Linux 上运行 \cmd{lscpu}，记录 socket、core、thread、cache 和 NUMA 信息。写一段短报告解释：如果 LCQI 的线性层使用 1、2、4、8 个线程，你预计最先遇到的瓶颈可能是什么。
\end{exercisebox}
